%% ============================================================
%% quasar/l1-l2-covalidation.tex
%% ============================================================
%%
%% Canonical L2 co-validation security theorem.
%% ============================================================

\subsection*{L2 co-validation security}

\begin{theorem}[L2 co-validation security inheritance]
\label{thm:l2-covalidation}
Let $\mathsf{Chain}_X$ be an L2 with
$\mathsf{primary\_network\_id} = N_P$, and let
$V_X \subseteq V_{N_P}^{(\eta_P, g_P)}$ at the moment of
$\mathsf{Chain}_X$'s key-era Bootstrap. If a Horizon-final certificate
on $\mathsf{Chain}_X$ binds
$\mathsf{primary\_nebula\_root} = R_P$ (a NebulaRoot of $N_P$ at era
$\eta_P$ generation $g_P$), then any adversary forging conflicting
finality on $\mathsf{Chain}_X$ must:
\begin{enumerate}[nosep]
  \item Compromise enough of $V_X$ to forge an L2-local Pulse, AND
  \item Either compromise enough of $V_{N_P}$ to forge an alternative
        $R_P$ (whence the security reduction inherits Theorem
        \ref{thm:horizon-soundness} applied to $N_P$), OR
  \item Break the transcript binding (collision on TupleHash256
        \cite{nistsp800185}).
\end{enumerate}
Therefore the L2's safety bound is at least the minimum of $N_P$'s
safety bound and the L2's own validator-set safety bound.

\textbf{Note.} This is co-validation security, not rollup security.
There is no fraud-proof window, no validity-proof anchoring; the
inheritance is purely via shared validator-set membership and
transcript binding.
\end{theorem}

\begin{remark}[L3 is a different layer model]
$\mathsf{Chain}_X$ at L3 (rollup / validity / app-execution layer per
Definition~\ref{def:network-tiers}) inherits security via a
\emph{succinct proof} of state transitions, not via co-validation
membership. L3 chains MAY use Groth16 or STARKs for compression. PQ
finality of an L3 inherits from the underlying L1/L2's Pulse + ML-DSA;
the proof system itself is not the PQ root of trust
(Remark~\ref{rem:groth16-wrapper}).
\end{remark}
